\documentclass[conference]{IEEEtran}
\IEEEoverridecommandlockouts
% La línea anterior solo es necesaria para identificar financiamiento en el
% primer pie de página. Si no aplica, se puede comentar.
\usepackage{cite}
\usepackage{amsmath,amssymb,amsfonts}
\usepackage{graphicx}
\usepackage{textcomp}
\usepackage{xcolor}
\def\BibTeX{{\rm B\kern-.05em{\sc i\kern-.025em b}\kern-.08em
    T\kern-.1667em\lower.7ex\hbox{E}\kern-.125emX}}
\begin{document}

\title{Percepción a Escala Real: un Simulador de Conducción en Realidad Virtual para el Aprendizaje del Tamaño y las Distancias del Vehículo}

\author{
\IEEEauthorblockN{1\textsuperscript{st} Integrante 1}
\IEEEauthorblockA{\textit{Pontificia Universidad Católica de Chile} \\
Santiago, Chile \\
correo1@uc.cl}
\and
\IEEEauthorblockN{2\textsuperscript{nd} Integrante 2}
\IEEEauthorblockA{\textit{Pontificia Universidad Católica de Chile} \\
Santiago, Chile \\
correo2@uc.cl}
\and
\IEEEauthorblockN{3\textsuperscript{rd} Integrante 3}
\IEEEauthorblockA{\textit{Pontificia Universidad Católica de Chile} \\
Santiago, Chile \\
correo3@uc.cl}
}

\maketitle

\begin{abstract}
Resumen breve del proyecto: problema, enfoque XR propuesto y aporte esperado.
\end{abstract}

\begin{IEEEkeywords}
realidad virtual, realidad extendida, simulación de conducción, percepción espacial, educación vial
\end{IEEEkeywords}

% ============================================================
% I. INTRODUCCIÓN (~0.5 pág.)
% ============================================================
\section{Introducción}
\label{sec:introduccion}

Aprender a conducir exige que la persona construya, además del conocimiento
de las reglas de tránsito, un sentido corporal del vehículo: cuánto mide,
dónde terminan sus esquinas, qué tan cerca puede pasar de un obstáculo o de
otro auto sin colisionar. Este componente propioceptivo suele adquirirse
recién dentro del vehículo real, bajo la supervisión de un instructor, lo
que concentra el aprendizaje en sesiones costosas, de disponibilidad
limitada y con margen de error nulo ante cualquier equivocación de un
conductor novato. Los simuladores de conducción en pantalla 2D (juegos de
manejo, simuladores de escritorio) permiten practicar reglas y reacciones,
pero al mostrarse en una pantalla plana y a una escala que no es 1:1,
no logran transmitir el tamaño real del vehículo ni la profundidad
necesaria para juzgar distancias con el cuerpo, que es precisamente la
habilidad más difícil de enseñar fuera del auto real.

Proponemos abordar este problema mediante un simulador de conducción en
realidad virtual dirigido a personas que están aprendiendo a conducir
(alumnos de escuelas de manejo, conductores primerizos), en el que la
persona experimenta el habitáculo y el volumen del vehículo a escala 1:1
mediante un visor de VR con seguimiento de cabeceo, en vez de una pantalla
2D. A diferencia de un simulador tradicional, la realidad virtual permite
que la persona perciba, con su propio cuerpo, dónde terminan las esquinas
del auto y cuán cerca pasa de un obstáculo — una noción espacial que un
monitor no puede transmitir porque no reproduce la estereopsis ni la
escala real del entorno. Con esto se espera reducir la brecha de
percepción espacial con la que los conductores novatos llegan a sus
primeras clases prácticas, disminuyendo el número de sesiones necesarias
con instructor y el riesgo asociado a los primeros errores de juicio de
distancia.

% ============================================================
% II. TRABAJOS RELACIONADOS Y MARCO TEÓRICO (~1 pág.)
% ============================================================
\section{Trabajos relacionados y marco teórico}
\label{sec:relacionados}

% TODO: citar con \cite{} y completar report/references.bib
% - Simuladores de conducción existentes (2D y VR) y sus limitaciones.
% - Literatura de percepción de profundidad y escala en VR (estereopsis,
%   presencia, propiocepción).
% - Brecha que aborda este proyecto: escala 1:1 + juicio corporal de
%   distancias, no solo reacción a reglas de tránsito.

% ============================================================
% III. PROPUESTA DE SOLUCIÓN (~2.5 pág.)
% ============================================================
\section{Propuesta de solución}
\label{sec:propuesta}

\subsection{Casos de uso}

\textbf{Caso de uso 1: Nombre.}
\begin{itemize}
    \item \textbf{Usuario:} \dots
    \item \textbf{Acción:} \dots
    \item \textbf{Objetivo:} \dots
\end{itemize}

\textbf{Caso de uso 2: Nombre.}
\begin{itemize}
    \item \textbf{Usuario:} \dots
    \item \textbf{Acción:} \dots
    \item \textbf{Objetivo:} \dots
\end{itemize}

\textbf{Caso de uso 3: Nombre.}
\begin{itemize}
    \item \textbf{Usuario:} \dots
    \item \textbf{Acción:} \dots
    \item \textbf{Objetivo:} \dots
\end{itemize}

\subsection{Descripción del sistema}
% Interacciones principales, hardware y software previstos,
% diagramas o wireframes, y un storyboard como figura.

% Ejemplo de figura (storyboard):
% \begin{figure}[htbp]
% \centerline{\includegraphics{figures/storyboard.png}}
% \caption{Storyboard de interacción típica del usuario con el sistema.}
% \label{fig:storyboard}
% \end{figure}

\subsection{Tecnologías y alcance}
% Justificación técnica del hardware/software (alternativas consideradas,
% criterios de selección, limitaciones) y qué queda dentro/fuera de
% alcance del prototipo.

% ============================================================
% IV. PLAN DE TRABAJO (~0.5 pág.)
% ============================================================
\section{Plan de trabajo}
\label{sec:plan}

% Ejemplo de tabla de hitos:
% \begin{table}[htbp]
% \caption{Cronograma hacia las próximas entregas}
% \begin{center}
% \begin{tabular}{|l|l|l|}
% \hline
% \textbf{Actividad} & \textbf{Responsable(s)} & \textbf{Fecha} \\
% \hline
% \dots & \dots & \dots \\
% \hline
% \end{tabular}
% \label{tab:cronograma}
% \end{center}
% \end{table}

% ============================================================
% REFERENCIAS (no cuentan para el límite de 5 páginas)
% ============================================================
\bibliographystyle{IEEEtran}
\bibliography{references}

\end{document}
